% Generated draft: requires mathematical and source review.
\documentclass[11pt]{article}
\usepackage[margin=25mm]{geometry}
\usepackage{fontspec}
\setmainfont{Latin Modern Roman}
\usepackage{amsmath,amssymb,mathtools,unicode-math}
\setmathfont{Latin Modern Math}
\usepackage{enumitem,needspace,xurl,hyperref}
\setlength{\parindent}{0pt}
\setlength{\parskip}{5pt}
\setlength{\emergencystretch}{3em}
\newcommand{\N}{\mathbb{N}}
\newcommand{\Z}{\mathbb{Z}}
\newcommand{\Q}{\mathbb{Q}}
\newcommand{\R}{\mathbb{R}}
\newcommand{\C}{\mathbb{C}}
\newcommand{\F}{\mathbb{F}}
\newcommand{\A}{\mathbb{A}}
\newcommand{\PP}{\mathbb{P}}
\newcommand{\E}{\mathbb{E}}
\newcommand{\eps}{\varepsilon}
\newcommand{\dd}{\,\mathrm{d}}
\begin{document}
\section*{UnsolvedMath ID: 6700051}
\textbf{Source problem code:} catalogue entry
\subsection{Definitions and mathematical statement}
Let $(S^n,g_{\mathrm{sph}})$ be the unit round sphere, so $\operatorname{Sc}_{g_{\mathrm{sph}}}=n(n-1)$, and let $B(R)$ be a geodesic ball of radius $R\ge\pi/2$. The source calls these concave spherical balls. The open-versus-closed ball convention and treatment of $R=\pi$ or larger are not stated; in the closed-ball convention $R\ge\pi$ gives the whole sphere.

The intended length comparison is usually the quadratic-form inequality
\[
 g_x(v,v)\ge(g_{\mathrm{sph}})_x(v,v)
 \qquad(x\in B(R),\ v\in T_xB(R)),
\]
although the source's word ``greater'' does not distinguish strict from non-strict comparison. Its length-extremality conjecture says that no such Riemannian metric can have
\[
 \operatorname{Sc}_g(x)>n(n-1)\qquad\text{at every }x\in B(R).
\]
Thus, under the non-strict metric comparison reading, the target is
\[
 \forall g\ge g_{\mathrm{sph}}|_{B(R)}\quad
 \exists x\in B(R):\operatorname{Sc}_g(x)\le n(n-1).
\]
The statement forbids strict scalar-curvature improvement everywhere; it does not assert a pointwise upper bound everywhere or an isometry conclusion in the equality case.

No agreement of metrics at the boundary, mean-curvature bound, prescribed boundary geometry, or extension across the boundary is supplied. The ball convention, smoothness up to any included boundary, dimension range, and metric-comparison convention must be fixed. These omissions are not filled by borrowing boundary hypotheses from a separate hemisphere rigidity theorem.
\subsection{Short English statement}
Can a metric that length-dominates the round metric on a spherical ball of radius at least $\pi/2$ have scalar curvature strictly above the spherical value everywhere? The conjecture says no, with no supplied additional boundary condition.

The research summary reports sphere and some hemisphere/ball results but gives no precise hypotheses, and says the general radius/dimension case remains open. It names Llarull, ``Brendle--Marques--Nevens,'' and ``Schoen--Yau/Huisken'' without exact references. Those informal rigidity attributions are not treated as verification of this boundary-unconstrained extremality statement. The partial-progress label and 2017 passage are retained with current status needing review.
\subsection{Sources}
\begin{itemize}
\item Local UnsolvedMath registry, permanent ID 6700051, source code catalogue entry, scalar-curvature Question [?54], ``Extremality of Concave Spherical Balls,'' proposed by Misha Gromov, 2017. Dataset curation and metadata: Ulam AI and UnsolvedMath Contributors, CC BY 4.0.
\item Misha Gromov, \emph{101 Questions, Problems and Conjectures around Scalar Curvature} (2017), printed tag [?54], occurrence 53, PDF page 57, \url{https://www.ihes.fr/~gromov/wp-content/uploads/2018/08/101-problemsOct1-2017.pdf}; normalized pdf-text extraction, accessed 2026-07-29. The extracted final reference marker is incomplete.
\item The summary's Llarull, ``Brendle--Marques--Nevens,'' and ``Schoen--Yau/Huisken'' attributions have no full entries or exact hypotheses. Spellings and asserted roles are not independently verified or corrected.
\item Source-specific provenance: the author-hosted PDF is incomplete and publicly available; current status and statement-reuse terms NEEDS\_REVIEW, separately from the dataset licence.
\end{itemize}
\end{document}
